\chapter{Penurunan subgradien terproyeksi}

% segment-id: d90.hab.v1.ch04.seg0001
Misalkan \(V\) ruang Hilbert real, \(C\subset V\) himpunan tak kosong, tertutup, dan konveks, serta \(f:V\rightarrow\R\) fungsi konveks. Kita tertarik menyelesaikan masalah
\begin{equation}
    \min_{x\in C} f(x).
\end{equation}
Kita mengasumsikan bahwa himpunan peminim tidak kosong dan menetapkan
\(x^*\in\arg\min_{x\in C}f(x)\). Masalah berkendala di atas selalu dapat ditulis sebagai masalah tak berkendala melalui
\begin{equation}
    \min_{x\in V} f(x) + \delta_C(x),
\end{equation}
dengan \(\delta_C\) fungsi indikator himpunan \(C\).\footnote{Formulasi ini memakai \(f\) yang terdefinisi pada seluruh \(V\), termasuk di luar \(C\).}

Ingat kembali metode optimisasi numerik paling dasar untuk masalah semacam ini ketika \(C=V\) dan \(f\) terdiferensialkan: penurunan gradien (atau penurunan \emph{tercuram}). Kita memilih nilai awal \(x_0\), lalu untuk \(k=0,1,\dots\) melakukan pembaruan
\begin{equation}
    x_{k+1} = x_k-\tau_k \nabla f(x_k),
\end{equation}
dengan \(\tau_k>0\) menyatakan ukuran langkah. Algoritma \emph{baru} pertama yang kita tinjau ialah versi subgradien dari penurunan tercuram: untuk \(k=0,1,\dots\),
\begin{equation}
    \begin{cases}
        g_k &\in \partial f(x_k),\\
        x_{k+1} &= x_k-\tau_k g_k.
    \end{cases}
\end{equation}
Jika \(C\neq V\), masalah segera dapat muncul: mungkin saja \(x_{k+1}\notin C\). Dalam hal itu, iterasi yang memakai subdiferensial \(f+\delta_C\) tidak lagi terdefinisi karena \(\partial\delta_C(x)=\emptyset\) untuk \(x\notin C\). Karena itu, kita perlu memastikan setiap iterasi baru tetap berada di \(C\); untuk tujuan ini kita memakai \emph{proyeksi}.

% segment-id: d90.hab.v1.ch04.seg0002
\begin{theorem}[Teorema proyeksi]
    Misalkan \(V\) ruang Hilbert real dan \(C\subset V\) tak kosong, tertutup, dan konveks. Fungsi berikut terdefinisi dengan baik:
    \begin{equation}
        \begin{aligned}
            \proj_C:V&\rightarrow C,\\
            x&\mapsto\arg\min_{y\in C}\|x-y\|.
        \end{aligned}
    \end{equation}
\end{theorem}
\begin{proof}
    Kita perlu menunjukkan bahwa minimum tercapai secara tunggal. Tetapkan
    \(d=\inf_{y\in C}\|y-x\|\), dan ambil barisan peminimum \((y_n)_n\subset C\), yaitu
    \begin{equation}
        \lim_n \|y_n-x\| = d.
    \end{equation}
    Karena \(C\) konveks, \((y_n+y_m)/2\in C\). Identitas jajaran genjang memberi
    \[
        \|y_n-y_m\|^2
        =2\|y_n-x\|^2+2\|y_m-x\|^2
        -4\left\|\frac{y_n+y_m}{2}-x\right\|^2
        \leq 2\|y_n-x\|^2+2\|y_m-x\|^2-4d^2.
    \]
    Ruas kanan menuju nol ketika \(m,n\rightarrow\infty\), sehingga \((y_n)_n\) merupakan barisan Cauchy. Kelengkapan \(V\) menghasilkan \(y_n\rightarrow\widehat y\in V\), sedangkan ketertutupan \(C\) menghasilkan \(\widehat y\in C\). Dari kontinuitas norma diperoleh
    \begin{equation}
        \|\widehat y-x\| = \lim_n \|y_n-x\| = d,
    \end{equation}
    jadi peminim memang ada.

    Untuk ketunggalan, andaikan \(y_1\neq y_2\) keduanya peminim. Identitas jajaran genjang memberikan
    \begin{equation}
        \begin{aligned}
            \left\|\frac{y_1+y_2}{2}-x\right\|^2
            &=\frac{1}{2}\|y_1-x\|^2+\frac{1}{2}\|y_2-x\|^2
              -\frac{1}{4}\|y_1-y_2\|^2\\
            &<d^2.
        \end{aligned}
    \end{equation}
    Namun \((y_1+y_2)/2\in C\), yang bertentangan dengan definisi \(d\). Jadi peminimnya tunggal.
\end{proof}

% segment-id: d90.hab.v1.ch04.seg0003
Dengan proyeksi, kita memperoleh cara mengatasi masalah \emph{keluar} dari \(C\): metode subgradien terproyeksi. Pilih \(x_0\in C\); kemudian, untuk \(k=0,1,\dots\) dan \(\tau_k\geq0\), tetapkan
\begin{equation}\label{subgradient:eq:projectedSG}
    \begin{cases}
        g_k &\in \partial f(x_k),\\
        x_{k+1} &= \proj_C( x_k-\tau_k g_k) .
    \end{cases}
\end{equation}
Untuk membuktikan konvergensi metode ini, kita memerlukan beberapa hasil bantu.

\begin{lemma}
    Misalkan \(V\) ruang Hilbert real dan \(C\subset V\) tak kosong, tertutup, dan konveks. Maka
    \begin{equation}
        \inner{x-\proj_C(x)}{y-\proj_C(x)}\leq 0,\quad y\in C.
    \end{equation}
    Khususnya, jika \(C\) subruang,\footnote{Artinya, \(C\) memenuhi definisi ruang vektor.} maka
    \begin{equation}
        \inner{x-\proj_C(x)}{y}= 0,\quad y\in C.
    \end{equation}
\end{lemma}
\begin{proof}
    Untuk setiap \(y\in C\), kita mempunyai
    \begin{equation}
        \begin{aligned}
            \|\proj_C(x)-x\|^2
            &\leq \|y-x\|^2\\
            &= \|y - \proj_C(x) + \proj_C(x)-x\|^2\\
            &= \|y-\proj_C(x)\|^2
              + 2 \inner{y-\proj_C(x)}{\proj_C(x)-x}
              + \|\proj_C(x)-x\|^2.
        \end{aligned}
    \end{equation}
    Akibatnya,
    \begin{equation}
        \begin{aligned}
            0\leq \|y-\proj_C(x)\|^2
            + 2 \inner{y-\proj_C(x)}{\proj_C(x)-x}.
        \end{aligned}
    \end{equation}
    Ambil sebarang \(z\in C\), lalu pada ketaksamaan di atas masukkan
    \(y=\proj_C(x) + t(z-\proj_C(x))\) untuk \(t\in (0,1)\). Kita memperoleh
    \begin{equation}
        \begin{aligned}
            0\leq t^2 \|z-\proj_C(x)\|^2
            + 2 t\inner{z-\proj_C(x)}{\proj_C(x)-x}.
        \end{aligned}
    \end{equation}
    Setelah dibagi dengan \(t>0\),
    \begin{equation}
        \begin{aligned}
            0\leq t \|z-\proj_C(x)\|^2
            + 2 \inner{z-\proj_C(x)}{\proj_C(x)-x}.
        \end{aligned}
    \end{equation}
    Dengan mengambil limit \(t\rightarrow 0\), kita memperoleh hasil yang diinginkan. Kesamaan untuk kasus subruang ditinggalkan sebagai latihan.
\end{proof}

% segment-id: d90.hab.v1.ch04.seg0004
\begin{prop}
    Misalkan \(V\) ruang Hilbert real dan \(C\subset V\) tak kosong, tertutup, dan konveks. Proyeksi pada \(C\) bersifat tak ekspansif, yaitu kontinu Lipschitz dengan konstanta satu:
    \begin{equation}
        \|\proj_C(x_1)-\proj_C(x_2)\|\leq \|x_1-x_2\|.
    \end{equation}
\end{prop}
\begin{proof}
    Tanpa mengurangi keumuman, andaikan
    \(\|\proj_C(x_2)-\proj_C(x_1)\|^2 \neq 0\), sebab jika tidak, tidak ada yang perlu dibuktikan. Kita memperoleh
    \begin{equation}
        \begin{aligned}
            \|\proj_C(x_2)-\proj_C(x_1)\|^2
            &= \inner{\proj_C(x_2)-\proj_C(x_1)}{\proj_C(x_2)-x_1+x_1-\proj_C(x_1)}\\
            &= \inner{\proj_C(x_2)-\proj_C(x_1)}{\proj_C(x_2)-x_1}\\
            &\quad+ \inner{\proj_C(x_2)-\proj_C(x_1)}{x_1-\proj_C(x_1)}\\
            &\leq \inner{\proj_C(x_2)-\proj_C(x_1)}{\proj_C(x_2)-x_1}\\
            &= \inner{\proj_C(x_2)-\proj_C(x_1)}{\proj_C(x_2)-x_2+x_2-x_1}\\
            &\leq \inner{\proj_C(x_1)-\proj_C(x_2)}{x_2-\proj_C(x_2)}\\
            &\quad+ \inner{\proj_C(x_2)-\proj_C(x_1)}{x_2-x_1}\\
            &\leq \inner{\proj_C(x_2)-\proj_C(x_1)}{x_2-x_1}\\
            &\leq \|\proj_C(x_2)-\proj_C(x_1)\|\|x_2-x_1\|.
        \end{aligned}
    \end{equation}
    Membagi dengan \(\|\proj_C(x_2)-\proj_C(x_1)\|\neq 0\) menghasilkan kesimpulan.
\end{proof}

% segment-id: d90.hab.v1.ch04.seg0005
\begin{lemma}[Ketaksamaan fundamental metode subgradien terproyeksi]\label{lemma:subgradient:fundamental}
    Misalkan \(x^*\in\arg\min_{x\in C}f(x)\), dan misalkan \(x_k\) adalah iterasi metode subgradien terproyeksi~\eqref{subgradient:eq:projectedSG}, dengan \(g_k\in\partial f(x_k)\) dan \(\tau_k\geq0\). Maka
    \begin{equation}
        \|x_{k+1}-x^*\|^2
        \leq \|x_k- x^*\|^2 - 2\tau_k (f(x_k)-f(x^*)) + \tau_k^2\|g_k\|^2.
    \end{equation}
\end{lemma}
\begin{proof}
    Dari definisi subgradien,
    \(f(x_k)+ \inner{g_k}{x^*-x_k}\leq f(x^*)\). Karena \(x^*\in C\), kita mempunyai \(\proj_C(x^*)=x^*\). Ketakspansifan proyeksi lalu memberikan
    \begin{equation}\label{eq:subgradient:fundamental}
        \begin{aligned}
            \|x_{k+1}-x^*\|^2
            &= \|\proj_C(x_k-\tau_kg_k) - \proj_C(x^*)\|^2\\
            &\leq \|x_k-\tau_k g_k - x^*\|^2\\
            &= \|x_k- x^*\|^2 - 2\tau_k \inner{g_k}{x_k-x^*} + \tau_k^2\|g_k\|^2\\
            &\leq \|x_k- x^*\|^2 - 2\tau_k (f(x_k)-f(x^*)) + \tau_k^2\|g_k\|^2.
        \end{aligned}
    \end{equation}
\end{proof}

% segment-id: d90.hab.v1.ch04.seg0006
Hal yang wajar dilakukan sekarang ialah memakai hasil sebelumnya untuk menentukan ukuran langkah optimal. Lebih tepatnya, dengan meminimumkan ruas kanan~\eqref{eq:subgradient:fundamental} terhadap \(\tau_k\), kita memperoleh
\begin{equation}
    \tau_k = \frac{f(x_k)-f(x^*)}{\|g_k\|^2}
\end{equation}
jika \(g_k\neq 0\), dan kita menetapkan \(\tau_k=0\) jika \(g_k=0\). Ukuran langkah ini disebut aturan ukuran langkah Polyak.

Sebagai latihan, buktikan bahwa jika fungsi konveks \(f:V\rightarrow\R\) bersifat \(L\)-Lipschitz secara global pada \(V\), maka setiap \(g\in\partial f(x)\) memenuhi \(\|g\|\leq L\). Pada hasil-hasil berikut, kita memakai langsung batas pada subgradien yang dipilih, sehingga juga mencakup situasi berkendala yang tidak memenuhi asumsi Lipschitz global tersebut.

\begin{theorem}
    Andaikan subgradien yang dipilih pada iterasi memenuhi \(\|g_k\|\leq L\). Dengan aturan ukuran langkah Polyak
    \begin{equation}
        \tau_k = \frac{f(x_k)-f(x^*)}{\|g_k\|^2}\quad\text{jika }g_k\neq0,
        \qquad \tau_k=0\quad\text{jika }g_k=0.
    \end{equation}
    metode subgradien terproyeksi~\eqref{subgradient:eq:projectedSG} memenuhi
    \begin{enumerate}
        \item \(\|x_{k+1}-x^*\|^2\leq \|x_k-x^*\|^2\), dengan ketaksamaan ketat jika \(f(x_k)>f(x^*)\);
        \item \(f(x_k)\rightarrow f(x^*)\) ketika \(k\rightarrow\infty\); dan
        \item \(f_{\mathrm{best}}^N-f(x^*)\leq \dfrac{L\|x_1-x^*\|}{\sqrt{N}}\), dengan \(f_{\mathrm{best}}^N = \min_{k=1,\dots,N}f(x_k)\).
    \end{enumerate}
\end{theorem}
\begin{proof}
    Tanpa mengurangi keumuman, kita boleh mengasumsikan \(f(x^*)=0\). Jika tidak, cukup tinjau \(\widetilde f(x)\coloneq f(x)-f(x^*)\). Dari~\cref{lemma:subgradient:fundamental} dan aturan ukuran langkah diperoleh
    \begin{equation}
        \begin{aligned}
            \|x_{k+1}-x^*\|^2
            &\leq \|x_k - x^*\|^2 - 2\tau_k f(x_k) +\tau_k^2\|g_k\|^2\\
            &= \|x_k - x^*\|^2 - a_k,
        \end{aligned}
    \end{equation}
    dengan \(a_k=f(x_k)^2/\|g_k\|^2\) jika \(g_k\neq0\), dan nol jika \(g_k=0\). Jika \(g_k=0\), kekonveksan memberi \(f(x_k)=f(x^*)\). Karena itu, pengurangan bersifat ketat tepat ketika celah nilai fungsi positif.

    Dengan menata ulang dan menjumlahkan untuk \(k=1,\dots,N\), kita memperoleh
    \begin{equation}
        \sum_{k=1}^{N} a_k
        \leq \|x_1-x^*\|^2 - \|x_{N+1}-x^*\|^2
        \leq \|x_1-x^*\|^2.
    \end{equation}
    Karena \(\|g_k\|\leq L\),
    \begin{equation}
        \sum_{k=1}^{N} f(x_k)^2
        \leq \sum_{k=1}^{N} L^2 a_k
        \leq L^2 \|x_1-x^*\|^2.
    \end{equation}
    Batas ini berlaku untuk setiap \(N\), sehingga \(f(x_k)\rightarrow0\). Selain itu,
    \begin{equation}
        N (f_{\mathrm{best}}^N)^2
        \leq \sum_{k=1}^{N} f(x_k)^2
        \leq L^2\|x_1-x^*\|^2,
    \end{equation}
    yaitu
    \begin{equation}
        f_{\mathrm{best}}^N\leq \frac{L\|x_1-x^*\|}{\sqrt{N}}.
    \end{equation}
\end{proof}

% segment-id: d90.hab.v1.ch04.seg0007
\begin{cor}
    Untuk mencapai ketelitian \(f_{\mathrm{best}}^N - f(x^*)\leq \epsilon\), kita memerlukan
    \[
        N=\Oc\!\left(\frac{L^2\|x_1-x^*\|^2}{\epsilon^2}\right)
    \]
    iterasi.
\end{cor}

\begin{rem}
    Ingat bahwa metode gradien mulus standar mencapai kompleksitas \(\Oc(1/\epsilon)\).
\end{rem}

Walaupun hasil konvergensi dengan aturan ukuran langkah Polyak menarik secara teoretis dan memberi batas kasus terburuk untuk iterasi terbaik, dalam praktiknya ukuran langkah tersebut tidak dapat dihitung tanpa mengetahui \(f(x^*)\). Berikut ini kita berikan hasil konvergensi yang lebih umum.

\begin{theorem}
    Andaikan subgradien yang dipilih memenuhi \(\|g_k\|\leq L\), dan ukuran langkah \(\tau_k>0\) memenuhi
    \[
        \frac{\sum_{k=1}^n\tau_k}{\sum_{k=1}^n\tau_k^2}\rightarrow \infty
        \qquad\text{ketika }n\rightarrow\infty.
    \]
    Maka
    \[
        f^n_{\mathrm{best}}-f(x^*)
        \leq
        \frac{\|x_1-x^*\|^2+L^2\sum_{k=1}^n\tau_k^2}
             {2\sum_{k=1}^n\tau_k}
        \longrightarrow0.
    \]
\end{theorem}
\begin{proof}
    Sekali lagi, untuk menyederhanakan notasi, andaikan \(f(x^*)=0\). Dengan memakai \(\|g_k\|\leq L\) dan menjumlahkan ketaksamaan fundamental~\cref{lemma:subgradient:fundamental} untuk \(k=1,\dots,n\), kita memperoleh
    \begin{equation}
        2\sum_{k=1}^n\tau_k f(x_k)
        \leq \|x_1-x^*\|^2 + L^2 \sum_{k=1}^n \tau_k^2.
    \end{equation}
    Tuliskan
    \[
        \sigma_n = \frac{\sum_{k=1}^n\tau_k}{\sum_{k=1}^n\tau_k^2}.
    \]
    Karena \(f^n_{\mathrm{best}}\leq f(x_k)\),
    \begin{equation}
        f^n_{\mathrm{best}}
        \leq \frac{\sum_{k=1}^n\tau_kf(x_k)}{\sum_{k=1}^n\tau_k}
        \leq \frac{\|x_1-x^*\|^2}{2\sum_{k=1}^n\tau_k}
        +\frac{L^2}{2\sigma_n}.
    \end{equation}
    Hipotesis \(\sigma_n\rightarrow\infty\) juga memaksa \(\sum_{k=1}^n\tau_k\rightarrow\infty\); karena itu, kedua suku pada ruas kanan menuju nol.
\end{proof}

% segment-id: d90.hab.v1.ch04.seg0008
Varian hasil konvergensi juga dapat diperoleh, misalnya jika \(C\) kompak; lihat~\cite{beck2017first} untuk perinciannya.

Konvergensi dapat dipercepat dan juga \emph{ditransfer} dari nilai fungsi ke iterasi \((x_k)_k\) jika kita menambahkan asumsi bahwa \(f\) konveks kuat.

\begin{theorem}
    Andaikan \(f\) konveks kuat dengan parameter \(\mu>0\) pada \(C\), dan subgradien yang dipilih memenuhi \(\|g_k\|\leq L\).\footnote{Batas ini dikenakan pada subgradien yang benar-benar dipilih oleh algoritma; batas tersebut tidak disimpulkan hanya dari sifat Lipschitz relatif pada \(C\).} Mulai dari \(x_1\in C\), gunakan ukuran langkah
    \(\tau_k = \dfrac{2}{\mu k}\) untuk \(k\geq1\). Untuk \(n\geq2\), metode subgradien terproyeksi~\eqref{subgradient:eq:projectedSG} memenuhi
    \begin{enumerate}
        \item \(f^n_{\mathrm{best}}-f(x^*)\leq \dfrac{2L^2}{\mu (n-1)}\); dan
        \item \(\|x_n-x^*\|\leq \dfrac{2L}{\mu \sqrt{n-1}}\).
    \end{enumerate}
\end{theorem}
\begin{proof}
    Seperti sebelumnya, tanpa mengurangi keumuman, tetapkan \(f(x^*)=0\). Kekonveksan kuat memperbaiki ketaksamaan fundamental menjadi
    \begin{equation}
        \begin{aligned}
            \|x_{k+1}-x^*\|^2
            &\leq (1-\mu\tau_k)\|x_k- x^*\|^2 - 2\tau_k f(x_k) + \tau_k^2\|g_k\|^2\\
            &\leq (1-\mu\tau_k)\|x_k- x^*\|^2 - 2\tau_k f(x_k) + \tau_k^2L^2.
        \end{aligned}
    \end{equation}
    Dengan menata ulang dan membagi dengan \(2\tau_k\), diperoleh
    \begin{equation}
        \begin{aligned}
            f(x_k)
            \leq -\frac{\tau_k^{-1}}{2}\|x_{k+1}-x^*\|^2
            + \frac{\tau_k^{-1}-\mu}{2}\|x_k- x^*\|^2
            + \frac{\tau_k}{2}L^2.
        \end{aligned}
    \end{equation}
    Memasukkan \(\tau_k = 2/(\mu k)\) menghasilkan
    \begin{equation}
        \begin{aligned}
            f(x_k)
            \leq -\frac{\mu k}{4} \|x_{k+1}-x^*\|^2
            + \frac{\mu(k-2)}{4}\|x_k- x^*\|^2
            + \frac{1}{\mu k}L^2.
        \end{aligned}
    \end{equation}
    Kalikan dengan \(k-1\), lalu jumlahkan terhadap \(k=1,\dots,n\). Suku jarak meneleskop, sehingga
    \begin{equation}
        \sum_{k=1}^{n} (k-1)f(x_k)
        \leq -\frac{\mu n(n-1)}{4} \|x_{n+1}-x^*\|^2
        + \frac{1}{\mu}\sum_{k=1}^{n}\frac{k-1}{k}L^2.
    \end{equation}
    Karena \(\sum_{k=1}^{n} (k-1)=n(n-1)/2\) dan
    \(\sum_{k=1}^{n}(k-1)/k\leq n\), kita memperoleh
    \begin{equation}\label{subgradient:eq:strongly}
        \frac{n(n-1)}{2} f^n_{\mathrm{best}}
        \leq \sum_{k=1}^{n} (k-1)f(x_k)
        \leq -\frac{\mu n(n-1)}{4} \|x_{n+1}-x^*\|^2
        + \frac{nL^2}{\mu}.
    \end{equation}
    Dengan membuang suku jarak yang tak positif, diperoleh
    \begin{equation}
        f^n_{\mathrm{best}} \leq \frac{2L^2}{\mu(n-1)}.
    \end{equation}

    Untuk batas iterasi, buang suku \(-2\tau_k f(x_k)\leq0\) dari ketaksamaan fundamental yang diperbaiki. Dengan \(d_k=\|x_k-x^*\|\),
    \begin{equation}
        d_{k+1}^2\leq\left(1-\frac{2}{k}\right)d_k^2
        +\frac{4L^2}{\mu^2k^2}.
    \end{equation}
    Untuk \(k=1\), ketaksamaan ini memberi \(d_2^2\leq4L^2/\mu^2\). Jika
    \(d_k^2\leq4L^2/[\mu^2(k-1)]\), substitusi pada relasi di atas memberikan
    \[
        d_{k+1}^2
        \leq \frac{4L^2}{\mu^2}
        \left(\frac{k-2}{k(k-1)}+\frac{1}{k^2}\right)
        \leq \frac{4L^2}{\mu^2k}.
    \]
    Induksi membuktikan \(\|x_n-x^*\|\leq2L/[\mu\sqrt{n-1}]\) untuk setiap \(n\geq2\).
\end{proof}
